\documentclass[11pt]{article}
\usepackage[utf8]{inputenc}
\usepackage{amsmath}
\usepackage{amssymb}
\usepackage{graphicx}
\usepackage{geometry}
\usepackage{hyperref}
\usepackage{enumitem}
\usepackage{booktabs}

\geometry{margin=1in}

\title{Force Sensing for Continuum Robots: \\
Analysis of Aloi's Method and Future Directions}
\author{Tongyu Wang\\Georgia Institute of Technology}
\date{\today}

\begin{document}

\maketitle

\section{Executive Summary}

This document analyzes the limitations of Aloi et al.'s Gaussian-based force estimation method for continuum robots and proposes a new approach that incorporates environmental knowledge to achieve unique and precise force estimation. The key insight is that while Aloi's method works well for certain scenarios, it fundamentally suffers from non-uniqueness issues when multiple solutions exist, 3D loads are present, or tangential friction forces occur.

\section{Analysis of Aloi's Method (2022)}

\subsection{Core Methodology}

Aloi's method estimates external loads on continuum robots using:
\begin{itemize}
    \item \textbf{Input}: Fiber Bragg Grating (FBG) sensor measurements of robot shape/curvature
    \item \textbf{Parameterization}: External loads as sum of Gaussian distributions
    \item \textbf{Optimization}: Levenberg-Marquardt algorithm to minimize shape prediction error
\end{itemize}

The load distribution is modeled as:
\begin{equation}
f(s) = \sum_{i=1}^{N} A_i \exp\left(-\frac{(s-\mu_i)^2}{2\sigma_i^2}\right)
\end{equation}

where $\theta = [A_1, \mu_1, \sigma_1, \ldots, A_N, \mu_N, \sigma_N]^T$ are the parameters to be estimated.

\subsection{Optimization Formulation}

The method solves:
\begin{equation}
\min_{\theta} \quad J(\theta) = \|\mathbf{y}_{\text{meas}} - \mathbf{y}_{\text{pred}}(\theta)\|^2 + \lambda \|\theta - \theta_0\|^2
\end{equation}

where:
\begin{itemize}
    \item $\mathbf{y}_{\text{meas}}$: measured robot shape from FBG sensors
    \item $\mathbf{y}_{\text{pred}}(\theta)$: predicted shape from Cosserat rod model with load $f(s)$
    \item $\lambda$: regularization weight
    \item $\theta_0$: prior estimate
\end{itemize}

\subsection{Strengths}

\begin{enumerate}
    \item \textbf{Smooth representation}: Gaussian parameterization provides continuous, differentiable load distribution
    \item \textbf{Computational efficiency}: Low-dimensional parameter space (6 parameters for 2 loads)
    \item \textbf{Good shape reconstruction}: Achieves excellent shape RMSE ($<0.15$ mm in our tests)
    \item \textbf{Handles distributed loads}: Natural representation for contact forces
\end{enumerate}

\subsection{Fundamental Limitations}

\subsubsection{1. Non-Uniqueness of Solutions}

\textbf{Problem}: Multiple load configurations can produce identical robot shapes.

\textbf{Mathematical Basis}: The inverse problem is ill-posed. Given shape measurements $\mathbf{y}$, the mapping $f(s) \rightarrow \mathbf{y}$ is not injective. Different load distributions can satisfy:
\begin{equation}
\mathbf{y}_{\text{pred}}(f_1) = \mathbf{y}_{\text{pred}}(f_2) \quad \text{but} \quad f_1 \neq f_2
\end{equation}

\textbf{Example Scenarios}:
\begin{itemize}
    \item Two small loads vs. one large load at intermediate position
    \item Distributed load vs. concentrated load with same moment
    \item Loads at different positions with compensating magnitudes
\end{itemize}

\textbf{Evidence from Our Implementation}:
\begin{itemize}
    \item Case 1: Load 2 magnitude error 91.7\% (before multi-start optimization)
    \item Case 3: Load 1 position error 40mm (13\% of robot length)
    \item Multi-start optimization (3600 initial guesses) required to find good solutions
\end{itemize}

\subsubsection{2. Point Load vs. Gaussian Mismatch}

\textbf{Problem}: Real contact forces are often point loads, but Gaussian representation is inherently distributed.

\textbf{Consequence}:
\begin{itemize}
    \item Cannot precisely localize point contacts
    \item Estimated load ``spreads out'' over $\pm 3\sigma$ region
    \item Position estimates are centroids, not actual contact points
\end{itemize}

\textbf{Our Results}: Even with optimal parameters, estimated loads show:
\begin{itemize}
    \item $\sigma \approx 10-20$ mm (3-7\% of robot length)
    \item Position errors up to 15mm for small loads
    \item Magnitude errors when loads are close together (Gaussian overlap)
\end{itemize}

\subsubsection{3. 2D Planar Assumption}

\textbf{Problem}: Method assumes all loads are normal to robot backbone in a single plane.

\textbf{Limitations}:
\begin{itemize}
    \item Cannot distinguish 3D out-of-plane loads
    \item Cannot estimate tangential (friction) forces along backbone
    \item Cannot handle torsional loads
\end{itemize}

\textbf{Real-World Impact}: Most manipulation tasks involve:
\begin{itemize}
    \item 3D contact forces (e.g., grasping, pushing objects)
    \item Friction during sliding contact
    \item Combined normal + tangential loads
\end{itemize}

\subsubsection{4. Local Minima in Optimization}

\textbf{Problem}: Non-convex objective function with many local minima.

\textbf{Manifestation}:
\begin{itemize}
    \item Sensitive to initial guess $\theta_0$
    \item Can converge to wrong load configuration
    \item Requires multi-start strategy (computationally expensive)
\end{itemize}

\textbf{Our Solution}: Implemented 3600 initial guesses (4 forces $\times$ 5 positions $\times$ 3 widths $\times$ 2 loads), but this:
\begin{itemize}
    \item Increases computation time 3600$\times$
    \item Still no guarantee of global optimum
    \item Not practical for real-time applications
\end{itemize}

\subsubsection{5. Number of Loads Must Be Known A Priori}

\textbf{Problem}: Must specify $N$ (number of Gaussians) before optimization.

\textbf{Consequence}:
\begin{itemize}
    \item If $N$ too small: cannot represent actual load distribution
    \item If $N$ too large: overfitting, spurious loads appear
    \item No automatic model selection mechanism
\end{itemize}

\section{Representative Failure Scenarios}

\subsection{Scenario 1: Multiple Solutions Exist}

\textbf{Setup}: Robot contacts a flat surface at unknown location.

\textbf{Problem}:
\begin{itemize}
    \item Single normal force at position $s_1$ produces shape $\mathbf{y}$
    \item But distributed load over $[s_1 - \Delta, s_1 + \Delta]$ can produce similar $\mathbf{y}$
    \item Aloi's method may converge to either solution
\end{itemize}

\textbf{Why It Fails}: No mechanism to prefer point contact over distributed contact.

\subsection{Scenario 2: 3D Loads}

\textbf{Setup}: Robot pushes against object at angle (force has $x$, $y$, $z$ components).

\textbf{Problem}:
\begin{itemize}
    \item FBG sensors measure bending in two planes: $\kappa_x(s)$, $\kappa_y(s)$
    \item But cannot measure torsion $\tau(s)$ or axial strain $\epsilon(s)$
    \item 3D force $\mathbf{f} = [f_x, f_y, f_z]^T$ projects to 2D measurements
\end{itemize}

\textbf{Why It Fails}: Underdetermined system - 3 unknowns, 2 measurements per point.

\subsection{Scenario 3: Tangential Friction}

\textbf{Setup}: Robot slides along surface with friction.

\textbf{Problem}:
\begin{itemize}
    \item Normal force $f_n(s)$ causes bending
    \item Tangential force $f_t(s)$ causes compression/tension
    \item FBG measures only curvature, not axial strain
\end{itemize}

\textbf{Why It Fails}: Tangential forces are invisible to curvature-based estimation.

\subsection{Scenario 4: Close Multiple Contacts}

\textbf{Setup}: Robot grasps object with two contact points 2cm apart.

\textbf{Problem}:
\begin{itemize}
    \item Two Gaussians with $\sigma \approx 1.5$ cm overlap significantly
    \item Optimization struggles to separate them
    \item Often converges to single larger Gaussian at midpoint
\end{itemize}

\textbf{Why It Fails}: Gaussian basis functions are not orthogonal; interference between nearby loads.

\section{Proposed Solution: Environment-Informed Force Estimation}

\subsection{Core Idea}

\textbf{Fundamental Insight}: Incorporate knowledge of the environment to constrain the solution space and achieve unique, precise force estimation.

\begin{equation}
\boxed{\text{FBG Measurements} + \text{Environment Knowledge} \rightarrow \text{Unique Force Estimate}}
\end{equation}

\subsection{Types of Environmental Knowledge}

\subsubsection{1. Geometric Constraints}

\textbf{Information}: Known object shapes, surface locations, workspace boundaries.

\textbf{Application}:
\begin{itemize}
    \item Contact can only occur at robot-object intersection points
    \item Force direction must be normal to contact surface
    \item Reduces search space from continuous $s \in [0, L]$ to discrete contact points
\end{itemize}

\textbf{Mathematical Formulation}:
\begin{equation}
s_i \in \mathcal{S}_{\text{contact}} = \{s : \|\mathbf{r}(s) - \mathbf{r}_{\text{obj}}(s)\| < \epsilon\}
\end{equation}

\subsubsection{2. Force Direction Constraints}

\textbf{Information}: Contact normal vectors from object geometry.

\textbf{Application}:
\begin{itemize}
    \item Force must be along surface normal: $\mathbf{f}_i = f_i \mathbf{n}_i$
    \item Reduces 3D force to 1D magnitude estimation
    \item Eliminates tangential component ambiguity
\end{itemize}

\textbf{Mathematical Formulation}:
\begin{equation}
\mathbf{f}(s_i) = f_i \mathbf{n}_{\text{obj}}(\mathbf{r}(s_i)) \quad \text{where} \quad f_i \geq 0
\end{equation}

\subsubsection{3. Contact Model Constraints}

\textbf{Information}: Physics-based contact models (Hertz, Hunt-Crossley, etc.).

\textbf{Application}:
\begin{itemize}
    \item Relate force magnitude to penetration depth: $f = k \delta^n$
    \item Constrain force-displacement relationship
    \item Provide prior on expected force magnitudes
\end{itemize}

\subsubsection{4. Friction Cone Constraints}

\textbf{Information}: Coulomb friction coefficient $\mu$.

\textbf{Application}:
\begin{itemize}
    \item Tangential force bounded: $|f_t| \leq \mu f_n$
    \item Enables estimation of both normal and tangential components
    \item Prevents physically impossible force configurations
\end{itemize}

\textbf{Mathematical Formulation}:
\begin{equation}
\|\mathbf{f}_t\| \leq \mu \|\mathbf{f}_n\| \quad \Rightarrow \quad \mathbf{f} \in \text{Friction Cone}
\end{equation}

\subsection{Proposed Method Formulation}

\subsubsection{Problem Statement}

Given:
\begin{itemize}
    \item FBG measurements: $\mathbf{y}_{\text{meas}} = [\kappa_x(s_1), \kappa_y(s_1), \ldots, \kappa_x(s_M), \kappa_y(s_M)]^T$
    \item Environment model: object geometry $\mathcal{O}$, contact model $\mathcal{C}$
    \item Robot model: Cosserat rod with parameters $(L, EI, GJ, EA)$
\end{itemize}

Estimate:
\begin{itemize}
    \item Number of contacts: $N$
    \item Contact locations: $\mathbf{s} = [s_1, \ldots, s_N]^T$
    \item Contact forces: $\mathbf{F} = [\mathbf{f}_1, \ldots, \mathbf{f}_N]^T$ (3D vectors)
\end{itemize}

\subsubsection{Optimization Formulation}

\begin{equation}
\begin{aligned}
\min_{N, \mathbf{s}, \mathbf{F}} \quad & J = \underbrace{\|\mathbf{y}_{\text{meas}} - \mathbf{y}_{\text{pred}}(\mathbf{s}, \mathbf{F})\|^2}_{\text{Shape fitting term}} + \underbrace{\lambda_1 \sum_{i=1}^{N} \|\mathbf{f}_i\|^2}_{\text{Force regularization}} \\
\text{subject to} \quad & s_i \in \mathcal{S}_{\text{contact}}(\mathcal{O}) \quad \forall i \\
& \mathbf{f}_i \cdot \mathbf{n}_i \geq 0 \quad \text{(no tension)} \\
& \|\mathbf{f}_{i,t}\| \leq \mu \|\mathbf{f}_{i,n}\| \quad \text{(friction cone)} \\
& f_i = \mathcal{C}(\delta_i) \quad \text{(contact model)} \\
& \delta_i = \text{penetration depth at } s_i
\end{aligned}
\end{equation}

\subsubsection{Key Advantages Over Aloi's Method}

\begin{enumerate}
    \item \textbf{Uniqueness}: Geometric constraints eliminate spurious solutions
    \item \textbf{Precision}: Point contact model (not Gaussian) for accurate localization
    \item \textbf{3D Capability}: Estimates full 3D force vectors
    \item \textbf{Friction Handling}: Explicitly models tangential forces
    \item \textbf{Automatic $N$ Selection}: Contact detection from geometry
\end{enumerate}

\subsection{Implementation Strategy}

\subsubsection{Phase 1: Contact Detection}

\begin{enumerate}
    \item Predict robot shape from FBG: $\mathbf{r}(s) = \int_0^s \mathbf{R}(s') \mathbf{e}_3 \, ds'$
    \item Find intersections with environment: $\mathcal{S}_{\text{contact}} = \{s : d(\mathbf{r}(s), \mathcal{O}) < \epsilon\}$
    \item Initialize $N = |\mathcal{S}_{\text{contact}}|$ contact points
\end{enumerate}

\subsubsection{Phase 2: Force Estimation}

For each contact point $s_i$:
\begin{enumerate}
    \item Compute surface normal: $\mathbf{n}_i = \nabla d(\mathbf{r}(s_i), \mathcal{O})$
    \item Compute penetration: $\delta_i = -d(\mathbf{r}(s_i), \mathcal{O})$
    \item Initialize force: $f_{i,n}^{(0)} = k \delta_i^{3/2}$ (Hertz model)
    \item Optimize: $\min_{\mathbf{f}_i} \|\mathbf{y}_{\text{meas}} - \mathbf{y}_{\text{pred}}\|^2$ subject to friction cone
\end{enumerate}

\subsubsection{Phase 3: Refinement}

\begin{enumerate}
    \item Jointly optimize all contact locations and forces
    \item Use gradient-based method (L-BFGS-B with box constraints)
    \item Validate solution: check force balance, moment balance
\end{enumerate}

\subsection{Required Inputs}

\begin{table}[h]
\centering
\begin{tabular}{@{}lll@{}}
\toprule
\textbf{Input} & \textbf{Source} & \textbf{Precision Required} \\ \midrule
Object geometry & CAD model / 3D scan & $\pm 1$ mm \\
Robot base pose & Robot kinematics & $\pm 2$ mm, $\pm 2^\circ$ \\
FBG measurements & Fiber sensors & $\pm 0.1$ mm (shape) \\
Material properties & Calibration & $\pm 10\%$ (EI, GJ) \\
Friction coefficient & Lookup table & $\pm 0.1$ \\ \bottomrule
\end{tabular}
\caption{Required environmental knowledge and precision}
\end{table}

\section{Validation Plan}

\subsection{Simulation Experiments}

\begin{enumerate}
    \item \textbf{Single contact}: Flat surface, varying angles
    \item \textbf{Multiple contacts}: Two-point grasp, three-point support
    \item \textbf{3D loads}: Out-of-plane forces, combined loading
    \item \textbf{Friction}: Sliding contact, tangential forces
    \item \textbf{Complex geometry}: Curved surfaces, corners, edges
\end{enumerate}

\subsection{Comparison Metrics}

\begin{itemize}
    \item \textbf{Position error}: $e_s = |s_{\text{est}} - s_{\text{true}}|$
    \item \textbf{Magnitude error}: $e_f = |\|\mathbf{f}_{\text{est}}\| - \|\mathbf{f}_{\text{true}}\||$
    \item \textbf{Direction error}: $e_\theta = \arccos(\hat{\mathbf{f}}_{\text{est}} \cdot \hat{\mathbf{f}}_{\text{true}})$
    \item \textbf{Shape RMSE}: $\text{RMSE} = \sqrt{\frac{1}{M}\sum_{j=1}^{M} \|\mathbf{r}_{\text{est}}(s_j) - \mathbf{r}_{\text{true}}(s_j)\|^2}$
\end{itemize}

\subsection{Baseline Comparisons}

\begin{itemize}
    \item Aloi's Gaussian method (current implementation)
    \item Rucker's EKF method
    \item Pure FBG shape sensing (no force estimation)
\end{itemize}

\section{Expected Outcomes}

\subsection{Quantitative Improvements}

\begin{table}[h]
\centering
\begin{tabular}{@{}lcc@{}}
\toprule
\textbf{Metric} & \textbf{Aloi Method} & \textbf{Proposed Method} \\ \midrule
Position error & 10-40 mm & $<$ 5 mm \\
Magnitude error & 10-90\% & $<$ 15\% \\
3D force estimation & No & Yes \\
Friction estimation & No & Yes \\
Uniqueness guarantee & No & Yes (with env.) \\ \bottomrule
\end{tabular}
\caption{Expected performance comparison}
\end{table}

\subsection{Qualitative Advantages}

\begin{itemize}
    \item \textbf{Physical interpretability}: Forces correspond to actual contacts
    \item \textbf{Robustness}: Fewer local minima due to constrained search space
    \item \textbf{Scalability}: Handles arbitrary number of contacts automatically
    \item \textbf{Real-time capability}: Faster convergence with good initialization
\end{itemize}

\section{Limitations and Future Work}

\subsection{Limitations of Proposed Method}

\begin{enumerate}
    \item \textbf{Requires environment model}: Not applicable to unknown environments
    \item \textbf{Sensitive to model errors}: Inaccurate geometry $\rightarrow$ wrong contact locations
    \item \textbf{Assumes rigid objects}: Deformable objects need different contact model
    \item \textbf{Computational cost}: Contact detection + optimization may be slow
\end{enumerate}

\subsection{Future Extensions}

\begin{enumerate}
    \item \textbf{Hybrid approach}: Use Aloi for unknown environments, proposed method when environment known
    \item \textbf{Online learning}: Update environment model from repeated contacts
    \item \textbf{Uncertainty quantification}: Bayesian framework for force estimation
    \item \textbf{Dynamic forces}: Extend to time-varying loads, impact forces
\end{enumerate}

\section{Conclusion}

Aloi's Gaussian-based method provides a strong foundation for force estimation but suffers from fundamental limitations:
\begin{itemize}
    \item Non-unique solutions
    \item Point load vs. Gaussian mismatch
    \item 2D planar assumption
    \item Cannot handle friction
\end{itemize}

The proposed environment-informed approach addresses these limitations by:
\begin{itemize}
    \item Constraining solution space with geometric knowledge
    \item Using point contact model for precision
    \item Estimating full 3D forces including friction
    \item Providing uniqueness guarantees
\end{itemize}

\textbf{Next Steps for Meeting}:
\begin{enumerate}
    \item Discuss which failure scenarios are most critical for target applications
    \item Determine availability of environment models (CAD, sensors, etc.)
    \item Define validation experiments and success criteria
    \item Assign implementation tasks and timeline
\end{enumerate}

\end{document}
